%======================================================================
%  Results -- Origin classification  (paper section 3; main figure = Figure 2)
%
%  Source artifact: Identifying plasmid origins with a genomic language model
%    https://claude.ai/code/artifact/aceb01ff-39f3-49c6-bfde-993d50131f6b
%    /home/canall/DNA_LM/BEND_ori_intervals/manuscript/artifact_origin/page.html
%    snapshot 2026-08-28 (page.html mtime 2026-08-28 15:33)
%
%  Seven paragraphs, no subsection headings, transcribed from the artifact's
%  "Main text" band. This artifact ships NO bibliography and NO inline
%  citation markers, so no \cite could be converted; every tool named below
%  carries a "% TODO cite:" at first mention. See CONVERSION_NOTES.md.
%
%  The artifact writes section pointers as "§2" and "§4". Under Nature-style
%  unnumbered sectioning there is no number to reference, so \secref resolves
%  them to the section name. TODO: the author may prefer different wording.
%
%  2026-09-01, author-directed: the seven transcribed paragraphs were
%  REGROUPED into four Results subsections, following the §2 split in the
%  same session:
%
%    3.1 Origin scoring in intergenic regions      sec:origin  (label kept)
%    3.2 Origins of transfer                       sec:origin-orit
%    3.3 Origins of replication                    sec:origin-oriv
%    3.4 Model scale, ensembling and inference cost  sec:origin-cost
%
%  Laid out 3 / 2 / 2 / 3 paragraphs (was 7; now 10). Two author-directed
%  goals drove the regrouping. First, the window -> intergenic-region
%  inference pipeline (ED Fig. 2) was described nowhere and its figure was
%  cited nowhere; it now opens the section, before any result, as 3.1 P2.
%  Second, oriT and oriV were reported together in every paragraph
%  ("0.87 on oriT and 0.77 on oriV"); they are now separated so that 3.2
%  carries the full baseline sweep and 3.3 states only what DIFFERS -- the
%  stronger baselines, the narrower margin, and the window/region ordering
%  -- without redefining the metrics, the splits or the novelty protocol.
%
%  Paragraph accounting. Old P1 -> 3.1 P1. Old P2 (two evaluation levels)
%  and old P5 s1-s2 (the novelty definition) merge into 3.1 P3. Old P3
%  (region level) splits oriT -> 3.2 P1, oriV -> 3.3 P1. Old P4 (window
%  level) splits oriT -> 3.2 P2, oriV -> 3.3 P2. Old P5's per-task numbers
%  follow their task. Old P6 splits: its scale/PlasmidGPT reading is pulled
%  forward into 3.4 P1, the ensemble into 3.4 P2, the timing into 3.4 P3.
%  Old P7 merges with old P6's last sentence, which said the same thing
%  twice ("this throughput difference was the reason we applied standalone
%  GPN-Plasmids across IMG/PR" / "We therefore applied the standalone
%  GPN-Plasmids classifiers across IMG/PR"). That is the ONLY transcribed
%  clause dropped, and it is dropped as a duplicate, not as content.
%
%  Added: 3.1 P2 in full (sourced from the ED Fig. 2 caption and from
%  supp_methods_03_origin.tex, "Calibration and region aggregation" and
%  "Inference" -- no fact that is not already in one of those); the
%  benchmark-design sentences opening 3.1 P3 (from the Fig. 2a caption and
%  Supplementary Methods, "Homology-partitioned cross-validation splits"),
%  including the oriT reinsertion caveat that previously appeared only in
%  Methods and in the Supp. Fig. S2 caption; the contrast sentences opening
%  3.3; and the two cross-cutting sentences opening 3.4 P1. One of those
%  last two is flagged inline for author sign-off.
%
%  Split by task, the per-task figures are tighter than the merged ranges
%  the transcription carried (old P4 gave "Evo2 probes scored 0.55--0.72"
%  and "one-hot probes 0.28--0.38" across BOTH tasks at once). Every value
%  below is read off tab:window-auprc / tab:region-map and their
%  supplementary counterparts at the two-decimal convention the Results
%  already use.
%
%  Eleven floats that no main-text sentence cited now have body citations:
%  ED Figs. 2-4 and Supp. Figs. S2-S8, plus Tables 1-2 and Supp. Tables
%  S2-S3. The citation order was chosen to reproduce the EXISTING float
%  numbering, so neither extended_data.tex nor supp_figures.tex needed its
%  floats reordered. See CONVERSION_NOTES.md.
%======================================================================


% ---- 3.1  Origin scoring in intergenic regions -----------------------
% \subsection*{A window-to-region pipeline scores plasmid intergenic regions for origins}
\subsection*{Origin scoring in intergenic regions}
\label{sec:origin}

Having pretrained \model on the diversity of IMG/PR, we next tested whether
its frozen representations could identify intergenic regions containing
origins of transfer or replication. Origins of transfer (\emph{oriT}) and
vegetative replication (\emph{oriV}) are the sequence elements that determine,
respectively, whether a plasmid can be mobilized between hosts and how it
replicates within one. They are short and intergenic, and outside a handful
of clinically important families they remain largely uncharacterized across
the metagenomic majority of plasmids. The task is to identify or prioritize
origin-containing intergenic regions rather than to recover exact origin
boundaries. We therefore trained lightweight probes on frozen embeddings,
without repeating the pretraining described in \secref{sec:pretraining}.
% NOTE (2026-09-01): the artifact read "the pretraining procedure of §2", and
% this line carried a TODO asking for better wording. Resolved when §2 was
% split into three subsections: \secref renders the heading TEXT, and
% "the pretraining procedure of Pretraining and model architecture" is a
% noun-phrase collision. Recast so the heading is the object of "in", which
% reads naturally for a title. sec:pretraining now labels §2.3.

A probe that scores fixed-length windows does not by itself annotate a
plasmid, because the unit of interest is a region rather than a window. We
therefore built the classifier into a two-stage inference pipeline
(\edfref{fig:ed-pipeline}). A query plasmid is first annotated into coding
sequences and intergenic regions. Each intergenic region is then tiled into
512 bp windows at a 16 bp stride and every window is scored from frozen
\model embeddings, with the forward and reverse-complement strands averaged
and each probe's logits mapped to probabilities by isotonic regression
(\nameref{sec:methods}; the task, training and inference parameters are
collected in \tref{tab:origin-params}). The result is a calibrated
probability profile along the region, which localizes the signal to part of
an intergenic region rather than only flagging the region as a whole. Pooling
that profile --- at its
75th percentile, or its maximum where a region has five or fewer windows ---
and recalibrating yields one score per intergenic region, which is what makes
the regions of a single plasmid comparable to each other and therefore
rankable. A plasmid carrying no characterized origin is thus reduced to a
short, ordered list of candidate regions, and this same procedure is the one
we later run across IMG/PR (\secref{sec:genomewide}).

We evaluated the pipeline on curated origins --- 112 \emph{oriT} and 224
\emph{oriV} after redundancy filtering --- partitioned so that held-out
origins are dissimilar to those used for training (\fref{fig:origin}a).
All-against-all local alignment between origin sequences assigned them to
five parts, and each cross-validation split trains on three parts while
validating on one and testing on one; because the parts are defined on the
origins rather than on the plasmids, a plasmid inherits its part through its
origin. For \emph{oriT}, 87 of the 112 origins partitioned under the
similarity threshold and the remaining 25 were reinserted next to their
nearest assigned neighbour, so that split is similarity-aware rather than
strictly homology-separated; all 224 \emph{oriV} origins partitioned without
reinsertion (\suppfrefs{fig:supp-clustermap-orit}{fig:supp-clustermap-oriv}).
Two complementary measurements follow from the pipeline's two stages.
Window-level detection asks whether origin-containing sequence can be
recovered among the intergenic windows of a plasmid, against a strongly
imbalanced background (held-out positive prevalence 0.048); it is summarized
by average precision (reported as AUPRC), for which a random ranker scores
the prevalence itself, 0.048 (\fref{fig:origin}b). Region-level ranking asks,
given a plasmid's intergenic regions, how well the origin-containing region
can be prioritized, and is scored by mean average precision (mAP) and by
Precision@1 (\fref{fig:origin}d). Because the parts were assigned by
origin-sequence similarity, 88\% (\emph{oriT}) and 87\% (\emph{oriV}) of
held-out positive windows had BLASTN
% TODO cite: BLASTN.
similarity $\leq$20 to their split's training positives, where similarity is
percent identity $\times$ query coverage / 100 (\fref{fig:origin}c). We call
these windows novel in this operational sense; the label does not establish
biological novelty. We report both measurements separately for each origin
type.


% ---- 3.2  Origins of transfer ----------------------------------------
% \subsection*{\model identifies origins of transfer where reference tools fail}
\subsection*{Origins of transfer}
\label{sec:origin-orit}

At the region level, a probe on frozen \model embeddings ranked the
\emph{oriT}-containing region first in most plasmids, reaching a Precision@1
of 0.87 across the five similarity-aware splits (\fref{fig:origin}d, top) and
a mean average precision of 0.90 (\tref{tab:region-map},
\supptref{tab:supp-region-metrics}). Existing annotation served this task
poorly. oriTfinder2
% TODO cite: oriTfinder2.
reached mAP 0.15, below the 0.35 of a naive baseline that simply ranks a
plasmid's longest intergenic region first --- on this benchmark the reference
tool is worse than a rule that reads nothing but region length. A probe
trained from
scratch on one-hot sequence reached 0.51 and a frozen PlasmidGPT
% TODO cite: PlasmidGPT.
probe 0.44. Among the pretrained nucleotide models \model matched or exceeded
Evo2
% TODO cite: Evo2.
probes up to 85$\times$ larger, with Evo2-7B and Evo2-40B both at 0.88.

The same ordering held at the window level, where \model reached AUPRC 0.78
on held-out test windows against the 0.048 random-ranker floor
(\fref{fig:origin}b, top; \tref{tab:window-auprc},
\supptref{tab:supp-window-auprc}); Evo2 probes scored 0.61--0.72 and one-hot
probes 0.28. PlasmidGPT was the weakest method at 0.15, falling below the
from-scratch one-hot baselines despite being a pretrained plasmid model. This
ranking holds split by split and not only on average
(\suppfref{fig:supp-prcurves}). Probes on the pretrained nucleotide models
also retained more of their performance on novel windows than the
from-scratch one-hot probes did: the lowest \model split on novel windows
scored 0.50, still above the best split of either one-hot probe (0.28), and
on novel windows \model remained similar to the larger Evo2 models (0.78
versus 0.77), within the observed split-to-split variation
(\fref{fig:origin}c, top). PlasmidGPT sat low on both axes. Because every
method's window probabilities are placed on a common calibrated scale before
aggregation
(\suppfref{fig:supp-reliability}), these probabilities can be read directly
along a region: at a held-out \emph{oriT} the profile rises over the
annotated origin and falls away across the flanking coding sequence
(\fref{fig:origin}e, top), as it does at four further held-out examples
(\edfref{fig:ed-tracks8}a--d).


% ---- 3.3  Origins of replication -------------------------------------
% \subsection*{Origins of replication are harder to rank against strong baselines}
\subsection*{Origins of replication}
\label{sec:origin-oriv}

The \emph{oriV} task has a different shape, because origins of replication
sit in intergenic regions that simpler methods already find fairly often.
Ranking a plasmid's longest intergenic region first reaches mAP 0.66, nearly
twice its \emph{oriT} value, and OriV-Finder
% TODO cite: OriV-Finder.
reaches 0.73, so unlike oriTfinder2 the reference tool here remains
comparatively competitive with the language models. Against that background
\model reached Precision@1 0.77 and mAP 0.82 (\fref{fig:origin}d, bottom;
\tref{tab:region-map}), ahead of the from-scratch one-hot probe (0.69) and
the PlasmidGPT probe (0.53). The \emph{oriT} result is therefore the stronger
of the two: on \emph{oriV} the margin over strong sequence-intrinsic
baselines is real but modest.

The two evaluation levels also separate the language models differently. At
the region level \model, Evo2-7B and Evo2-40B are indistinguishable by mAP
--- 0.816, 0.825 and 0.810, differences well inside the split-to-split spread
of any of them (\supptref{tab:supp-region-metrics}) --- whereas at the window
level \model leads by a clearer margin, AUPRC 0.67 against 0.62 for Evo2-7B,
with the other Evo2 scales lower still at 0.57 and 0.55, one-hot probes at
0.38 and PlasmidGPT again last at 0.20 (\fref{fig:origin}b, bottom;
\tref{tab:window-auprc}).
Robustness to novel windows followed the \emph{oriT} pattern: the lowest
\model split scored 0.58 against 0.45 for the best one-hot split, and on
novel windows \model matched the larger Evo2 models at 0.66
(\fref{fig:origin}c, bottom). Recall@1 and mAP for every method are reported
in \suppfref{fig:supp-regforest}, and ranking accuracy is broadly stable
across plasmid-length bins (\suppfref{fig:supp-maplen}). Calibrated
probability profiles again localize the origin within its region
(\fref{fig:origin}e, bottom; \edfref{fig:ed-tracks8}e--h).


% ---- 3.4  Model scale, ensembling and inference cost ------------------
% \subsection*{Ensembling helps, but standalone inference is 45$\times$ cheaper}
\subsection*{Model scale, ensembling and inference cost}
\label{sec:origin-cost}

Two features of the comparison hold across both origin types.
% FLAG (2026-09-01): the next sentence is the one genuinely NEW claim in this
% section. It is read directly off tab:region-map and tab:window-auprc
% (Evo2-40B <= Evo2-7B at both levels on both tasks) but the source artifact
% never stated it. Author sign-off needed; delete the sentence if unwanted.
Parameter count did not translate monotonically into accuracy, in that
Evo2-40B, the largest model tested, did not exceed Evo2-7B at either level on
either task (\tref{tab:region-map}, \tref{tab:window-auprc}). And pretraining
helped only where the pretraining corpus matched the target: PlasmidGPT is
pretrained on plasmids, but on engineered constructs rather than
environmental ones, and its probe ranked below a from-scratch one-hot
convolution on both tasks and at both levels. Neither result is an artefact
of how the probes were configured, since every embedder was given the layer
that maximized its own validation AUPRC, selected on validation columns only
(\edfref{fig:ed-layer}); nor of how the reference tools were run, since
restricting their databases to the training fold rather than to training plus
validation changed their scores by at most 0.007 (\suppfref{fig:supp-reftool}).

Calibrated probabilities from Evo2-7B and \model were combined by
arithmetic-mean soft voting. The ensemble exceeded both constituents at the
region level (mAP 0.93 for \emph{oriT} versus 0.90 for \model and 0.88 for
Evo2-7B; mAP 0.86 for \emph{oriV} versus approximately 0.82 for each
constituent; \tref{tab:region-map}) and at the window level (AUPRC 0.814 and
0.692 versus 0.781/0.671 and 0.722/0.616; \tref{tab:window-auprc}). This
improvement is consistent with partly non-redundant predictive information,
but it does not establish a particular mechanistic decomposition. One
hypothesis is corpus complementarity: \model was pretrained on plasmids with
diversity-aware sampling that preserves within-PTU variation, whereas Evo's
published OpenGenome construction retained one IMG/PR representative per PTU.
It is also possible, though untested, that Evo's additional pretraining on
whole prokaryotic genomes exposed it to origin-like sequence in integrated
plasmids or other chromosomal mobile elements. A direct test would stratify
ensemble gain by the presence of related chromosomally integrated elements or
other proxies for corpus-specific coverage.

The ensemble is a benchmark result; the genome-wide predictions used
downstream were generated with standalone \model. That choice was practical.
Evo2-7B has approximately 15 times as many parameters as \model (7.0B versus
468M), and a single-split, five-run window predictor scored a region in
49.8 ms with \model versus 2,231.7 ms with Evo2-7B (approximately
45$\times$; \fref{fig:origin}f), measured on 794 regions from 10 plasmids on
a CUDA device. Under the available compute budget that throughput difference
decided the deployment: we applied the standalone \model classifiers across
IMG/PR to generate genome-wide origin predictions, which form the basis for
the mobilization and replication analyses in \secref{sec:genomewide}.
% NOTE (2026-09-01): the artifact read "the mobilization and replication
% analyses in §4" and this line carried a wording TODO. No rewording was
% needed -- \secref already renders the heading as the object of "in", which
% is the construction the §2 split settled on. TODO downgraded to NOTE.
